We developed an end-to-end system for predicting the destination of a Porto taxi from a partially observed trajectory.  Starting from 1,710,670 labelled trips, we retained 1,672,901 after validity, spatial, and duration filtering.  We generated synthetic training prefixes with the same length distribution as the 320 supplied competition trajectories, then represented each prefix with its first and last GPS points, temporal variables, and dispatch metadata.  A shared 541-feature matrix supports six model families, while CatBoost uses a 36-column representation with native categorical variables.

We compared Linear Regression, Ridge, KNN, Decision Tree, Random Forest, XGBoost, and CatBoost under a common validation design, then combined the two strongest complementary models.  Random Forest achieved the best standalone results: 2.119~km on validation and 2.206~km on the 160-trip local public subset.  A fixed blend of 81.399\% Random Forest and 18.601\% XGBoost reduced the grouped cross-fitted validation error by 6.531~m on average and achieved the best local public score, 2.191~km.  The gain is consistent but small, so the main conclusion is that Random Forest provides the strongest individual model and XGBoost adds a modest corrective signal.  The principal limitations are the random rather than temporal validation split, repeated model selection on one holdout, and mixed CPU/GPU runtime measurements.
